\chapter{Inside the Compiler}\label{ch:internals}
\sloppy

This chapter describes how the compiler to C works: the stages a program
passes through, how Dylan values look in memory, how functions are
called and how tail calls are kept proper, how generic functions
dispatch, how exits, cleanups and conditions are built, and how the
garbage collector and threads fit in. You do not need any of it to use
the compiler. It helps when you want to know why one way of writing a
program is faster than another (Chapter~\ref{ch:performance}), when you
read the C the compiler writes, or when you want to work on the compiler
itself. The design, phase by phase, is also recorded in
\file{CompilerPlan.md} at the top of the source tree.

\section{The stages}\label{sec:internals-pipeline}

\index{compiler!stages}
The compiler is about 9\,000 lines of Common Lisp in \file{compiler/}.
It is a \term{whole-program compiler}: it sees the entire program, with
the library, before it writes any C, and since nothing changes once the
program runs, everything it learns about the program is true for good.
A build passes through these stages:

\begin{enumerate}
\item \textbf{Reading} (\file{front/source.lisp}). Each file is read
  with the reader of the interactive implementation
  (\file{src/compat.lisp}), the only code the two implementations share
  besides the macros' run time. The reader records the file, line and
  column where every list begins, so that every later diagnostic, and
  every \code{\#line} directive, can name a place in the source. A file's
  \code{module:} header (design note~33) says which module its forms
  belong to.

\item \textbf{Modules} (\file{front/modules.lisp}). Every
  \code{define-module} form in the program is found first, and the
  modules are defined, each after the modules it uses. Then every name
  in the program is resolved to a \emph{module variable}: one object per
  variable, owned by one module, however many modules import it and under
  whatever names. From here on names are gone, and nothing can be
  confused by renaming or shadowing.

\item \textbf{Macros and core forms} (\file{front/convert.lisp},
  \file{front/macros.lisp}). Each top-level form is converted to
  \emph{core Dylan}, about a dozen kinds of node: constants, variable
  references and assignments, \code{if}, sequences, \code{bind} (with
  multiple values), mutually recursive local methods, methods, calls,
  calls of runtime primitives, \code{bind-exit} and
  \code{unwind-protect}, and the definitions. Every other special form
  and built-in macro --- \code{cond}, \code{case}, \code{for},
  \code{for-each}, \code{handler-case}, backquote, \code{with-lock} and
  the rest --- is rewritten into these. A loop becomes a local method
  that calls itself in tail position. The program's own macros are
  expanded here too, by running their bodies in the interactive
  implementation (Chapter~\ref{ch:extensions}).

\item \textbf{The IR} (\file{ir/anf.lisp}, \file{ir/ir.lisp}). Core
  Dylan is converted to an intermediate representation in
  \term{A-normal form}: every intermediate value is given a name, and
  every operand of a call or primitive is a constant or a variable.
  Every variable is bound exactly once (the IR is \emph{alpha-converted}),
  so a variable is a unique object that analyses can annotate. Operands
  are evaluated left to right, as in the source.

\item \textbf{Analysis} (\file{ir/analyze.lisp}). For each variable:
  whether it is assigned, whether a function other than the one binding
  it uses it (it is \emph{captured}), and how often it is referred to.
  For each local function: where it is called, and whether it
  \emph{escapes} --- is used other than by being called. For each exit
  procedure of a \code{bind-exit}: whether it escapes its function. And
  for each call: whether it is in \emph{tail position}, that is, whether
  its value is the value of its function.

\item \textbf{Control lowering} (\file{ir/control.lisp}). A
  \code{bind-exit} whose exit is only called directly within its own
  function stays, and will become a jump. Every other \code{bind-exit},
  and every \code{unwind-protect}, is rewritten as a call of a runtime
  function, with the body moved into a new closure
  (Section~\ref{sec:internals-exits}).

\item \textbf{Contification} (\file{ir/contify.lisp}). A local function
  that never escapes and whose every call is a jump --- a call in tail
  position from the function that defines it, or from another such
  function --- is not compiled as a function at all, but as a labelled
  block of the function that contains it, its calls becoming
  assignments to its parameters and a \code{goto}. This is how loops
  become C loops.

\item \textbf{Type inference} (\file{backend/types.lisp},
  Section~\ref{sec:internals-types}) finds a class for as many
  variables as it can, and decides which \code{<double-float>} variables
  can be plain C \code{double}s.

\item \textbf{Calls} (\file{backend/calls.lisp},
  Section~\ref{sec:internals-calls}). Each call whose function is known
  is resolved: to a method, to a generic function's dispatch function,
  to a runtime primitive, or to inline C.

\item \textbf{Escape analysis} (in \file{backend/codegen.lisp}) finds
  \code{\#rest} lists, closures and instances that cannot outlive the
  function that makes them, and puts them on the C stack instead of the
  heap.

\item \textbf{Code generation} (\file{backend/codegen.lisp},
  \file{backend/objects.lisp}, \file{backend/cdata.lisp}). Every
  remaining function becomes a C function; every class, symbol and
  constant becomes static C data; every generic function gets a dispatch
  function.

\item \textbf{Output} (\file{backend/output.lisp}). The C is split into
  \file{\meta{name}.c}, one file per module, and \file{\meta{name}.h}
  (Section~\ref{sec:dylanc-files}), and \file{driver.lisp} runs the C
  compiler.
\end{enumerate}

When the library is compiled with the program (\cmd{-Os},
\cmd{-{}-whole}), tree shaking (\file{backend/shake.lisp},
Section~\ref{sec:internals-shaking}) runs between control lowering and
code generation.

The library itself is Dylan, in module \code{dylan}, in the files of
\file{runtime/dylan/} read in name order: conditions, numbers,
characters, collections, sequences, lists, vectors, strings, stretchy
vectors, deques, ranges, tables, arrays, printing and threads. It is
compiled with every program, by the same stages, on top of the
\code{\%} primitives written in C in \file{runtime/c/library.c} and the
other files of \file{runtime/c/}.

\section{How values are represented}\label{sec:internals-representation}

\index{representation of values}\index{tagged words}
Every Dylan value is one 64-bit machine word, the C type \code{D}
(\file{runtime/c/dylan.h}). The low bits of the word say what it is:

\begin{center}
\begin{tabular}{ll}
\toprule
Low bits & Value \\
\midrule
\code{...xx1} & an \code{<integer>}: the value shifted left one bit, so
  63 bits signed \\
\code{...010} & a \code{<character>}: the Unicode code point shifted left
  three bits \\
\code{...000} & a pointer to an object, 8-byte aligned \\
\bottomrule
\end{tabular}
\end{center}

Integers and characters are \term{immediate}: they need no memory of
their own, and allocating one costs nothing. The integer $n$ is the word
$2n+1$; that is why a debugger shows the integer 100 as 201. Adding two
tagged integers $2a+1$ and $2b+1$ is one machine addition and a
subtraction of 1, with the overflow checked by the C compiler's
\code{\_\_builtin\_add\_overflow}; an overflow signals an error rather
than wrapping.

Every object a pointer points to begins with a \term{header} word: a
pointer to its class's \term{wrapper}. The wrapper holds what the
runtime needs to know about the class --- the class object, the class's
numeric id, the instance size, how instances are represented (an ordinary
instance, a pair, a byte string, a double float, and so on), and the set
of the class's ancestors. Pairs are a header and two words; strings a
header, a length and the bytes; a \code{<double-float>} a header and a C
\code{double}; an instance of a class the program defines a header and
one word per slot.

\code{\#t}, \code{\#f} and \code{()} are three static objects, compared
by address. \code{\#t} and \code{\#f} are instances of an internal class
\code{<boolean>}, which the manual does not name and the \code{dylan}
module does not export.

\paragraph{Class ids.} \index{class ids}
Since the program is static, the compiler knows every class. It numbers
them in a preorder walk of the tree that each class makes with its
\emph{first} superclass, so that a class's descendants in that tree have
the ids from its own up to a limit. Where every subclass of a class lies
in that range, \code{(instance? x C)} is a range check on \code{x}'s
class id: one subtraction and one comparison. Where multiple inheritance
puts a subclass outside it, the test reads a bit in the ancestor set
every wrapper carries. The classes, their wrappers and the symbols are
all static C data emitted by the compiler, so a program does no work at
startup to build its class hierarchy.

\paragraph{Memory.}
Objects are allocated by the garbage collector
(Section~\ref{sec:internals-gc}), through the functions
\code{dy\_alloc} and its kin in \file{dylan.h}, which hide which
collector it is: those that can hold pointers apart from those that
cannot --- strings, float boxes, ratios --- so that the collector does
not scan the latter.

\section{The calling convention}\label{sec:internals-calls-convention}

\index{calling convention}\index{uniform entry}
Every Dylan function has a \term{uniform entry}, a C function of one
prototype:

\begin{ccode}
D f(D self, D a0, D a1, D a2, D a3, intptr_t argc);
\end{ccode}

\code{self} is the function object being called --- a closure, a
generic function, or a method --- through which a closure reaches its
free variables. The first four arguments travel in \code{a0} to
\code{a3}, in machine registers; \code{argc} says how many there are.
Arguments after the fourth go in a per-thread \term{argument area}, which
the callee reads before it calls anything else. Every function checks
\code{argc} on entry, so calling a function with the wrong number of
arguments is an error, not a crash.

One prototype for every function is what makes proper tail calls
possible (Section~\ref{sec:internals-tail}). It costs a little: a
function of two arguments is still passed six. So a function that is
called from places the compiler can see, other than in tail position,
also gets a \term{specialized entry}, named with \code{\_spec}, with
exactly its own parameters: \code{double}s where the arguments are
unboxed floats, and one parameter per keyword, filled in when compiling
from the keywords the call passes (\code{DY\_UNSUPPLIED} for those it
does not). A call through a specialized entry passes no count, checks
no arity and parses no keywords when it runs.

\paragraph{Multiple values.}
\index{multiple values}
A function returns its first value as its C return value. When it
returns some other number of values, it stores the count and all the
values in a per-thread values buffer. A caller that wants all the values
sets the count to 1 before the call; a function that returns one value
leaves the buffer alone unless it has itself made a call that may have
changed the count. A caller that wants only the first value --- the
common case --- reads the C return value and ignores the buffer.

\paragraph{\code{next-method}.}
A method that calls \code{next-method} receives its next method through
the thread: the dispatch function (or the method before it) stores it
there, and the method reads it on entry. The chain of methods for each
combination of argument classes is known when compiling, so the
next-method objects are static data, one for each tail of each chain.
Methods that do not use \code{next-method} ignore it, and a method
whose \code{next-method} is only called, with no arguments, reuses its
own arguments rather than capturing them in a closure.

\section{Proper tail calls}\label{sec:internals-tail}

\index{tail calls}\index{musttail}
Dylan requires that a call in tail position not use stack: a loop
written as a recursion must run in constant space, however many times it
goes round. C promises nothing of the kind. The compiler meets the
requirement in three ways.

\begin{enumerate}
\item A function's tail call of itself, and calls among local methods
  that have been contified, become \code{goto}s within one C function.
  This covers almost every loop: \code{for}, \code{while},
  \code{dotimes} and \code{for-each} all expand into local methods
  calling themselves in tail position.
\item Every other tail call --- through a module variable, a closure
  passed as an argument, a generic function, \code{next-method},
  \code{apply} --- is a C tail call marked
  \code{\_\_attribute\_\_((musttail))}, which GCC~15 and later and
  Clang~13 and later \emph{guarantee} to compile as a jump. The
  attribute requires caller and callee to have the same prototype, which
  the uniform entry provides. The generated C writes it with the macro
  \code{DY\_TAIL\_CALL}.
\item With \cmd{-{}-trampoline}, a tail call instead stores the function
  and its arguments in the thread and returns a marker; the nearest
  caller that is not making a tail call runs a loop that makes the
  pending call until a real value comes back. Any C11 compiler can
  compile this.
\end{enumerate}

A specialized entry cannot make a \code{musttail} call, because its
prototype differs from its callee's. When a function makes tail calls,
either its uniform entry holds the only body, or its specialized entry
gets a second copy of the body whose tail calls are ordinary calls. That
is safe: a specialized entry is only ever entered by a call that is not
a tail call, so each holds at most one extra C frame.

Calls inside an \code{unwind-protect}, inside a \code{bind-exit} whose
exit escapes, and inside a handler's extent are not in tail position:
their frame must remain until they return, as the language's semantics
already imply.

\section{Closures}\label{sec:internals-closures}

\index{closures}
Closures are \emph{flat}: a closure is a function object holding its C
entry and a copy of the value of each free variable it uses. A variable
that is both captured by a closure and assigned lives instead in a
one-word heap \term{cell}, which every closure that uses it shares, so
all of them see each assignment. A function with no free variables ---
every method defined at top level, for instance --- is a static object
and costs nothing to make. A local method that is only called, and never
escapes, is not made into a closure at all: it is contified or called
directly.

\section{Generic dispatch}\label{sec:internals-dispatch}

\index{generic functions!dispatch}\index{dispatch}
Because the program is static, the compiler knows every method of every
generic function. For each generic function it computes, at compile
time, which methods apply to each combination of argument classes and
singletons, sorts each such set by specificity, and writes a C
\term{dispatch function}: a decision tree that tests the arguments ---
singletons first, by comparing pointers, then a \code{switch} on the
class id --- and ends at a leaf for each distinct sorted list of
methods. A leaf checks keyword arguments, if the methods take any, sets
the thread's next method, and tail-calls the first method. A call that
no method fits reaches \code{dy\_no\_applicable\_method}, which signals
the error.

Here is the dispatch function for a generic function \code{area} with
one method on \code{<square>} and one on \code{<circle>}:

\begin{ccode}
D dy_gf_area_51a90f52_dispatch(D self, D a0, D a1, D a2, D a3,
                               intptr_t argc) {
    (void)a0; (void)a1; (void)a2; (void)a3;
    if (DY_UNLIKELY(argc != 1)) dy_wrong_argument_count(self, argc);
    D x0 = a0; (void)x0;
    switch (dy_wrapper_of(x0)->id) {
    case 61:
        goto leaf_0;
    case 62:
        goto leaf_1;
    default: goto no_method;
    }
leaf_0:
    dy_thread->next_method = DY_FALSE;
    DY_TAIL_CALL_DIRECT(fn1197_area, DY_OBJ(&dy_method_326),
                        a0, a1, a2, a3, argc);
leaf_1:
    dy_thread->next_method = DY_FALSE;
    DY_TAIL_CALL_DIRECT(fn1199_area, DY_OBJ(&dy_method_327),
                        a0, a1, a2, a3, argc);
no_method:
    dy_no_applicable_method(self);
}
\end{ccode}

\noindent (Three long lines have been broken to fit the page.)

There are no caches and no locks: the tree is fixed, and threads share
it freely.

Methods are ordered by the rule of the manual's chapter~10: arguments
compared left to right, a singleton before any class, a class before its
superclasses. The manual leaves open the order of two classes neither of
which is a subclass of the other; this compiler, like CLOS and like the
interactive implementation, puts first the one that comes earlier in the
\emph{argument's} class precedence list. Two methods that remain
unordered are ambiguous, and calling them is an error.

Better still is not to dispatch at all. Where type inference shows that
only one method can apply, the call goes straight to that method, or the
method is done inline (Section~\ref{sec:internals-calls}). A generic
function with a single method, specialized on \code{<object>} for every
argument, is never dispatched.

\paragraph{Classes and \code{make}.}
For each class that is instantiated, the compiler writes a
\term{maker}, which allocates the instance, fills each slot from its
init-keyword or its default, checks the init-keywords, and calls
\code{initialize} only if the program defines methods on it. A
\code{make} of a class the program defines, where only the library's
\code{make} method on \code{<class>} applies, calls the maker directly.
Slot offsets are fixed where every possible class of the receiver agrees
on them; otherwise accessors look the offset up in a small table in the
wrapper.

\section{Exits, cleanups and conditions}\label{sec:internals-exits}

\index{bind-exit@\code{bind-exit}!implementation}
\index{unwind-protect@\code{unwind-protect}!implementation}
A \code{bind-exit} whose exit procedure is only called directly, within
the function that binds it (or within local methods contified into it),
compiles to a jump: calling the exit assigns the values and does a
\code{goto} past the body. This is the common case --- an early return
from a loop --- and it costs nothing.

Every other \code{bind-exit}, and every \code{unwind-protect}, needs a
\term{frame} in the thread's \term{dynamic chain}, the list of active
exits, cleanups and handlers, innermost first. The frames live on the C
stack of three runtime functions in \file{runtime/c/control.c}:
\code{dy\_call\_with\_exit}, \code{dy\_unwind\_protect} and
\code{dy\_handler\_bind}. The compiler moves the body into a closure and
passes it to one of these:

\begin{dylan}
(unwind-protect p c)  =>  (%primitive unwind-protect (method () p)
                                                     (method () c))
(bind-exit (x) b)     =>  (%primitive call-with-exit (method (x) b))
\end{dylan}

\code{dy\_call\_with\_exit} pushes an exit frame, calls \code{\_setjmp},
and calls the body with an exit procedure. Calling that exit procedure
stores its arguments as the values to return, runs the cleanups of every
frame between it and the target, innermost first, and \code{longjmp}s to
the target frame. An exit called after its \code{bind-exit} has returned
is detected --- each frame carries a serial number --- and signalled as
an error. Because only these runtime functions call \code{setjmp},
generated code never does, and can make tail calls everywhere.

\paragraph{Conditions.}
\index{conditions!implementation}
Handlers are frames in the same chain. \code{signal} walks the chain
for a handler whose type matches and whose test, if any, is true, and
calls it with the condition and a \code{next-handler} that continues the
search further out; if none applies it calls \code{default-handler}. The
condition classes and \code{signal}, \code{error}, \code{cerror},
\code{check-type}, \code{abort}, \code{break} and
\code{default-handler} are written in Dylan, in
\file{runtime/dylan/01-conditions.dyl}; only the frame operations and
the transfer of control are C. When the runtime detects an error of its
own --- a wrong argument count, a type error, an overflow --- it calls
back into that Dylan code, so the error is an ordinary
\code{<simple-error>} or \code{<type-error>} that handlers can catch.

\section{The garbage collector}\label{sec:internals-gc}

\index{garbage collection}\index{Memory Pool System}
Memory is managed by the Memory Pool System in \file{runtime/mps/},
compiled as one C file into the runtime library, which is built with
\code{DY\_MPS}. It is a moving collector, so it must know every object's
layout: \file{runtime/c/gc-mps.c} gives it an \emph{object format} ---
for each kind of object, its size, the words of it that are references,
and how to leave a forwarding pointer or padding where an object was ---
and the \emph{roots}: the program's static data (found from the linker's
symbols on Linux, and from the Mach-O segments on macOS) and each
thread's stack and registers, scanned conservatively, and each thread's
state and thread variables, allocated outside the heap and registered
as a root. A tagged integer or character on a stack that happens to look
like a pointer into the heap can only keep an object alive, and in
place, longer than necessary; it can do no harm.

Objects go in five \emph{pools}: AMC, generational and moving, for
objects that may hold references; AMCZ for those that do not (strings,
floats, ratios), which the collector never scans; AMS, which does not
move, for \code{<thread>} objects, since a thread's state points to its
own; LO, which is neither moved nor protected, for objects holding a
mutex or condition variable the kernel uses; and AWL, weak, for the
table of identity hashes. Each thread has its own \emph{allocation
points} in AMC and AMCZ, from which it allocates by bumping a pointer;
the slow path, when the point's buffer is empty, enters the collector.
An object's fields are filled in before it is committed to the
collector, which may move it or ask for it to be allocated again.
Because the collector moves objects, a Dylan object's identity hash
cannot be its address: \code{\%object-hash} gives each object a number
the first time it is asked, kept in a table with weak keys, rebuilt when
the collector reports that keys have moved. The collector's behavior can
be tuned with environment variables (Chapter~\ref{ch:performance}).

\paragraph{The Boehm-Demers-Weiser collector.} \index{Boehm collector}
With \cmd{dylanc -{}-gc boehm} (Section~\ref{sec:dylanc-gc}) the
runtime is built without \code{DY\_MPS} and memory is managed instead by
\file{libgc}, built for threads. It is a \emph{conservative} collector:
it treats any word on a stack, in a register or in a scanned object that
looks like a pointer into the heap as one, so it needs no object format.
Strings and float boxes are allocated atomic so that their contents are
not mistaken for pointers. Each thread's state --- the values buffer, the
argument area, the dynamic chain --- is allocated uncollectable, so the
collector scans it and sees pointers held there. The collector does
mark-and-sweep collection, in parallel on two threads, and allocates
small objects from thread-local free lists without locking.

The functions \code{dy\_alloc} and \code{dy\_alloc\_init} in
\file{dylan.h} allocate with either collector.

\section{Threads}\label{sec:internals-threads}

\index{threads!implementation}
Each thread registers itself and its stack with the collector as it
starts, and makes its allocation points (under the Boehm collector,
\file{dylan.h} defines \code{GC\_THREADS} and includes \file{gc.h}
before \file{pthread.h}, so \code{pthread\_create} goes through the
collector, which registers each thread's stack). Each thread has one
\code{\_Thread\_local} pointer, \code{dy\_thread}, to its state: the
values buffer, the argument area, the dynamic chain, its next method,
its pending tail call in trampoline mode, its \code{<thread>} object,
and its values of the program's thread variables.

A thread starts by setting up that state and an exit frame at the top of
its dynamic chain, then calls its function. An unhandled serious
condition reaches \code{default-handler}, which in a thread other than
the main one records the condition in the \code{<thread>} and unwinds to
that top frame, ending the thread; \code{join-thread}, written in Dylan
in \file{runtime/dylan/20-threads.dyl}, finds it there and signals a
\code{<thread-error>}. Threads are detached when made; \code{join-thread}
waits on a condition variable in the \code{<thread>}. A \code{<lock>}
is an error-checking POSIX mutex, a \code{<recursive-lock>} a recursive
one, and an \code{<atomic-integer>} a C11 \code{\_Atomic} integer.

\section{Separate compilation of the library}\label{sec:internals-separate}

\index{separate compilation}\index{library cache}
Compiling the library's megabyte of C with every program would double
the time of a small build. So, by default, the compiler compiles the
library \emph{as if it did not know the program}. Every function the
\code{dylan} module exports, or that the compiler's expansions name, may
be called from anywhere with anything: no argument types and no
specialized entries are derived from the program's calls, and nothing of
the library is dropped. The C names of the library's constants,
variables and generic functions are derived from what they are, not
from the order in which the compiler met them, and its C temporaries are
numbered within each function. \file{\meta{name}-dylan.c} declares only
what it refers to. The result is that the library's C is the same text
for almost every program, and the driver compiles it once per C compiler
and option set and caches the object (Section~\ref{sec:dylanc-cache}).

What is lost is specialization of library code to the program, and the
dropping of unused library code. Calls \emph{from} the program
\emph{into} the library lose little, since the compiler still inlines
the library's arithmetic, \code{element}, \code{size} and iteration on
known collections at the call site.

\section{Tree shaking}\label{sec:internals-shaking}

\index{tree shaking}
With \cmd{-Os} or \cmd{-{}-whole} the library is compiled with the
program, and \file{backend/shake.lisp} keeps only what the program can
reach. The program's own top-level forms are all kept, as are the
library's top-level forms that are not definitions and every
\code{define-class} (classes are cheap static data). A library
definition of a variable, generic function or method is kept when
something kept refers to its variable. A method is kept, besides, only
if each class its specializers name could have instances --- a class
the runtime itself makes instances of, or one that something kept uses
as a value (makes, say, or answers). Naming a class as a specializer or
a superclass does not count.

Shaking is by name and class, not by data flow, and so it keeps more
than a program needs: a call of a generic function such as
\code{print-object}, \code{make} or \code{as} keeps every method that
could apply to a class with instances. That is why hello world is still
about 220~KB when built \cmd{-Os}.

\section{Type inference and inlining}\label{sec:internals-types}
\label{sec:internals-calls}

\index{type inference}
The type of a variable, for the compiler, is a class all its values are
instances of, or nothing known. Types come from:

\begin{itemize}
\item constants;
\item the result classes of the operations done inline --- an integer
  addition answers an integer, a float multiplication a
  \code{<double-float>};
\item \code{make} of a class whose maker the call goes to;
\item the getters of slots declared with a \code{type:};
\item the type a \code{bind} declares, as in
  \code{(bind (((x <integer>) \meta{init})) ...)}, which is checked when
  the binding is made;
\item the specializers of a method defined with \code{define-method},
  which dispatch has already checked;
\item for local functions, global functions and sole methods whose every
  call is known --- a loop, most often --- the classes of the arguments
  at those calls, and the classes of what the functions return.
\end{itemize}

These are found together, starting from nothing, until none changes.
The types are also \emph{flow-sensitive} where a test proves more: in
the branch an \code{(instance? x C)} test chooses, \code{x} is known to
be a \code{C}.

Calls whose function is known are then resolved (\file{calls.lisp}). A
call of \code{+}, \code{-}, \code{*}, \code{/}, \code{<}, \code{>} or
\code{=} with two arguments calls the binary generic function
(\code{binary+} and so on) directly. The library's methods for integer
and float arithmetic and comparison, for \code{size}, \code{element}
and the iteration protocol on lists, simple object vectors and byte
strings, and \code{id?}, \code{not}, \code{car}, \code{cdr},
\code{instance?} and slot getters and setters are done \emph{inline} ---
when the static program shows that no other method could apply to such
arguments. The manual's generic functions that the runtime writes in C
--- \code{instance?}, \code{object-class}, \code{subclass?} and the
other reflective functions of \dylanbook{chapters 10 and 11} --- each
have one library method, which calls a \code{\%} primitive; while the
program adds no method of its own to one, its calls are calls of the
primitive, and \code{instance?} of a known class is done inline. Each inline case is guarded by a test of the arguments' classes,
and the test is left out where their types are known. A float
operation with one argument known to be a \code{<double-float>} answers a
\code{<double-float>} (float contagion), so its result's type is known
even if the other argument's is not.

\paragraph{Unboxing.} \index{unboxing}
A \code{<double-float>} variable that no closure holds, and every use of
which wants a C \code{double}, is \emph{unboxed}: it is a C
\code{double} local, and computing it allocates nothing. A float
variable that is returned, stored in a slot or passed to an unknown
function needs its box, and is allocated on the heap each time it
changes. Integers need no unboxing; they are immediate.

\paragraph{Escape analysis.} \index{escape analysis}
A \code{\#rest} list that is only walked, through variables no closure
holds, is built on the C stack (up to eight pairs); keyword arguments
are parsed where they were passed, with no list at all. A closure, or an
instance whose class's maker \code{make} calls (and which has no
\code{initialize} methods), is made on the stack of the function that
makes it when nothing can reach it after that function returns: it is
only tested, called, given to a slot accessor sure to apply, or passed to
a known function whose parameter is used the same way. Constants
propagate through \code{bind}s, integer arithmetic on constants is done
when compiling, an \code{if} on a constant keeps one branch, and unused
pure bindings are dropped.

In a debug build (\cmd{-g}) none of the inferred types is used, so every
inline operation tests its arguments.

\section{Reading the generated C}\label{sec:internals-reading}

\index{C files!reading}
\cmd{dylanc -k} keeps the C. Here is a small program, \file{sum.dyl}:

\begin{dylan}
(define-method sum-to ((n <integer>))
  (for ((i 0 (+ i 1))
        (s 0 (+ s i)))
       ((> i n) s)))
(print (sum-to 100))
\end{dylan}

\noindent and the C function the compiler writes for \code{sum-to}, in
\file{sum-dylan-user.c}:

\begin{ccode}
D fn1190_sum_to(D self, D a0, D a1, D a2, D a3, intptr_t argc) {
    (void)self; (void)a0; (void)a1; (void)a2; (void)a3; (void)argc;
    D v105609, v105610, v105612, v105613, v105614, v105615, v105616,
      v105617, v105618, v105619;
    D t1;
    D t2;
    D t3;
    D t4;
#line 1 "sum.dyl"
    if (DY_UNLIKELY(argc != 1)) dy_wrong_argument_count(self, argc);
    v105609 = a0;
#line 2 "sum.dyl"
    t1 = DY_INT(0);
    t2 = DY_INT(0);
    v105612 = t1;
    v105613 = t2;
    goto for_1;
    for_1:;
#line 4 "sum.dyl"
    v105615 = DY_BOOL((intptr_t)(v105609) < (intptr_t)(v105612));
#line 2 "sum.dyl"
    if (v105615 != DY_FALSE) {
        dy_thread->mv_count = 1;
        return v105613;
    } else {
        v105617 = dy_int_add(v105612, DY_INT(1));
#line 3 "sum.dyl"
        v105619 = dy_int_add(v105613, v105612);
#line 2 "sum.dyl"
        t3 = v105617;
        t4 = v105619;
        v105612 = t3;
        v105613 = t4;
        goto for_1;
    }
}
\end{ccode}

\noindent (one line of declarations has been broken to fit the page).
Reading it from the top:

\begin{itemize}
\item The function has the uniform entry. Its name, \code{fn1190\_sum\_to},
  is a number and the Dylan name; every variable \code{v}\meta{n} is
  one IR variable, and every \code{t}\meta{n} a temporary.
\item It checks that it was given one argument; the method is the only
  method of \code{sum-to}, and its dispatch function has already checked
  that the argument is an \code{<integer>}.
\item The \code{for} loop has become a label, \code{for\_1}, and a
  \code{goto}. The step forms are evaluated into temporaries before any
  loop variable is assigned, because \code{for} steps its variables in
  parallel.
\item \code{(> i n)} is a single C comparison of the tagged words, with
  no test of their classes: \code{n} is an integer by its specializer,
  and \code{i} by type inference, since it starts at 0 and is only ever
  incremented. \code{(> i n)} became \code{(< n i)}, and on tagged
  integers $2a+1 < 2b+1$ exactly when $a<b$.
\item \code{(+ i 1)} and \code{(+ s i)} are \code{dy\_int\_add},
  an inline function in \file{dylan.h} that adds the tagged words and
  signals an error on overflow.
\item Returning one value sets the values count to 1 first, because a
  value count might be left over from a call made earlier in the
  function.
\item The \code{\#line} directives name the Dylan line each statement came
  from.
\end{itemize}

If \code{sum-to} were instead written as \code{(define sum-to (method
((n <integer>)) ...))}, \code{n} would not be known to be an integer
inside it: a method made by \code{method} and stored in a variable is
called through the variable, and must check its own argument. The
comparison would then be preceded by a test, \code{if
(DY\_IS\_INT(v105609))}, with a call of the generic function
\code{binary<} for the other case. Small differences like this are the
subject of the next chapter.

\fussy
